\chapter{支持向量机的间隔分类}
\label{chap:svm}
\label{chap:kernel-algorithms}
\begin{quote}\small
\textit{本章摘要。}当多条直线都能分开训练样本时，支持向量机用最近样本到边界的距离选择分类器。本章从几何间隔推导硬间隔优化，再通过松弛变量处理重叠类别，完整建立拉格朗日对偶、支持向量与核化预测之间的联系。
\end{quote}

上一章建立了核与特征空间内积的关系。本章研究二分类：给定输入$\bm x_i\in\R^d$和标签$y_i\in\{-1,+1\}$，怎样学习一个边界，将新的观测分到两类中？例如，输入是已经标准化的温度与振动，标签表示随后规定时间内是否发生故障。支持向量机（Support Vector Machine，SVM）的推导参照Theodoridis第11.10节的最优间隔分类及第11.11节的计算讨论\cite{theodoridis2015ml}，也可对照MML第12章\cite{deisenroth2020mml}。

\section{分类分数与几何间隔}
先用线性分数
\begin{equation}
 f(\bm x)=\bm w^{\mathsf T}\bm x+b,\qquad
 \hat y=\operatorname{sign} f(\bm x).
\end{equation}
其中$\bm w\in\R^d$是法向量，$b\in\R$是截距。$f=0$是决策超平面，$f>0$对应正类。分数等于零时需约定分类规则。若把$\bm w,b$同时乘以任意正数，预测类别不变，但分数会改变。因此不能直接比较不同参数下的分数大小来判断边界离样本多远。

设$\bm x_0$在边界上。从$\bm x$沿法向量移动到边界，写成$\bm x_0=\bm x-t\bm w$。代入边界方程得
\begin{equation}
 0=\bm w^{\mathsf T}\bm x+b-t\|\bm w\|_2^2,
 \qquad t=\frac{f(\bm x)}{\|\bm w\|_2^2}.
\end{equation}
于是点到超平面的距离为$|t|\|\bm w\|_2=|f(\bm x)|/\|\bm w\|_2$。对带标签的点，用
\begin{equation}
 \rho_i=y_if(\bm x_i),\qquad
 \gamma_i=\frac{y_if(\bm x_i)}{\|\bm w\|_2}
 \label{eq:svm-geometric-margin}
\end{equation}
分别表示函数间隔和有符号几何间隔，假定$\bm w\ne\bm0$。正确分类时两者为正，错分时为负。几何间隔不随共同正比例缩放而改变。

对严格可分的数据，可把参数缩放到$\min_i\rho_i=1$。此时最近点到边界的距离为$1/\|\bm w\|_2$，两侧间隔边界$f=\pm1$的距离为$2/\|\bm w\|_2$。图\ref{fig:svm-margin-geometry}给出一个具体设置。最大化间隔希望边界留出较大的几何余量，但实际泛化还取决于数据分布、标签质量和核的选择。

\begin{figure}[htbp]
\centering
\begin{tikzpicture}[scale=0.85,>=Latex]
 \draw[->] (-3,0)--(3.4,0) node[right] {$x_1$};
 \draw[->] (0,-1.8)--(0,2.5) node[above] {$x_2$};
 \draw[thick] (0,-1.5)--(0,1.9);
 \draw[dashed] (-1,-1.5)--(-1,1.9);
 \draw[dashed] (1,-1.5)--(1,1.9);
 \foreach \x/\y in {-1/1,-1/-1,-2/0.5} {\draw[blue,thick] (\x,\y) circle (0.09);}
 \foreach \x/\y in {1/1,1/-1,2/0.5} {\fill[red] (\x,\y) rectangle +(0.12,0.12);}
 \draw[<->] (-1,-1.35)--(1,-1.35) node[midway,below] {$2/\|\bm w\|_2=2$};
 \node[above] at (-1,1.9) {$f=-1$};
 \node[above] at (1,1.9) {$f=1$};
 \node[right] at (0,2.1) {$f=0$};
\end{tikzpicture}
\caption{教学数据的间隔几何。圆点为负类，方点为正类。取$\bm w=(1,0)^{\mathsf T},b=0$时，边界为$x_1=0$，两条虚线为$x_1=\pm1$，四个最近点位于间隔边界上。}
\label{fig:svm-margin-geometry}
\end{figure}

\section{硬间隔目标与对偶问题}
\subsection{从距离得到优化约束}
在两类均有样本、训练数据严格线性可分时，硬间隔SVM求解
\begin{equation}
 \min_{\bm w,b}\frac12\|\bm w\|_2^2,
 \qquad y_i(\bm w^{\mathsf T}\bm x_i+b)\ge1\quad(i=1,\ldots,n).
 \label{eq:svm-hard-margin-primal}
\end{equation}
平方范数与范数有相同的最小点，前面的$1/2$便于求导。约束固定了函数间隔尺度，因此范数最小对应几何间隔最大。若数据不可分，可行集合为空，不能继续用这个目标求分类器。

\subsection{样本展开与拉格朗日对偶}
为每个约束引入乘子$\alpha_i\ge0$，拉格朗日函数为
\begin{equation}
 \mathcal L=\frac12\|\bm w\|_2^2
 +\sum_i\alpha_i\{1-y_i(\bm w^{\mathsf T}\bm x_i+b)\}.
\end{equation}
对$\bm w,b$求偏导并令其为零，得到
\begin{equation}
 \bm w=\sum_i\alpha_i y_i\bm x_i,
 \qquad \sum_i\alpha_i y_i=0.
 \label{eq:svm-stationarity}
\end{equation}
若第二个条件不满足，对$b$取下确界会得到$-\infty$。代回第一式，利用$\sum_i\alpha_i y_i\bm w^{\mathsf T}\bm x_i=\|\bm w\|_2^2$，得到对偶函数
\begin{equation}
 g(\bm\alpha)=\sum_i\alpha_i-
 \frac12\sum_{i,j}\alpha_i\alpha_j y_i y_j\bm x_i^{\mathsf T}\bm x_j.
\end{equation}
对偶问题是最大化$g$，约束为$\alpha_i\ge0$及$\bm\alpha^{\mathsf T}\bm y=0$。原问题是凸二次规划。严格可分时存在参数使$m=\min_i y_i f_i>0$。将$\bm w,b$共同乘以大于$1/m$的数后，所有约束都严格满足，满足Slater条件，所以强对偶成立，原问题和对偶问题有相同最优值。

卡鲁什--库恩--塔克（Karush--Kuhn--Tucker，KKT）条件包含驻点条件、原始可行性、对偶可行性和互补松弛。这里的互补松弛是
\begin{equation}
 \alpha_i\{y_if(\bm x_i)-1\}=0.
 \label{eq:svm-hard-kkt}
\end{equation}
若$y_if(\bm x_i)>1$，必有$\alpha_i=0$。只有约束活跃的点可能有非零系数。称$\alpha_i>0$的点为支持向量。位于间隔边界上的点在退化情形下也可能系数为零，不能只看几何位置判断其对偶系数。

\subsection{两点算例}
一维训练集为$(x_1,y_1)=(-1,-1)$、$(x_2,y_2)=(1,+1)$。约束是$w-b\ge1,w+b\ge1$，因此$w\ge1$，最优解为$w=1,b=0$。对偶约束令$\alpha_1=\alpha_2=a$，目标变成$2a-2a^2$。令导数$2-4a=0$，得到$a=1/2$。由式\eqref{eq:svm-stationarity}重建$w=1$。两点都是支持向量，间隔宽度为2。

\section{软间隔目标与支持向量}
\subsection{松弛变量与合页损失}
引入$\xi_i\ge0$允许违反规范化间隔，软间隔目标为
\begin{equation}
 \min_{\bm w,b,\bm\xi}\frac12\|\bm w\|_2^2+C\sum_i\xi_i,
 \quad y_if(\bm x_i)\ge1-\xi_i,\quad\xi_i\ge0,
 \label{eq:svm-soft-margin-primal}
\end{equation}
其中$C>0$。固定$\bm w,b$时，最小可行松弛量是$\xi_i=\max(0,1-y_if(\bm x_i))$。消去$\xi_i$便得到合页损失（hinge loss）的正则化目标。

当$y_if\ge1$时损失为零。当$0<y_if<1$时分类正确，但位于间隔内部，损失介于0与1之间。当$y_if<0$时分类错误，损失大于1。$y_if=0$对应边界上的点，损失为1。松弛量总和衡量违约程度，不能把它当成错分个数。

若采用平均损失目标$\lambda\|\bm w\|^2/2+n^{-1}\sum_i\max(0,1-y_if_i)$，与式\eqref{eq:svm-soft-margin-primal}等价的参数关系为$C=1/(n\lambda)$。比较不同软件或样本量时要核对这一约定。

\subsection{软间隔对偶与箱约束}
用$\alpha_i\ge0$处理间隔约束，用$\beta_i\ge0$处理$\xi_i\ge0$：
\begin{equation}
 \mathcal L=\tfrac12\|\bm w\|^2+C\sum_i\xi_i
 +\sum_i\alpha_i(1-\xi_i-y_if_i)-\sum_i\beta_i\xi_i.
\end{equation}
对$\xi_i$求导得$C-\alpha_i-\beta_i=0$。结合非负性得到$0\le\alpha_i\le C$。对$\bm w,b$的驻点条件仍为式\eqref{eq:svm-stationarity}。因此
\begin{equation}
 \max_{\bm\alpha}\ \bm1^{\mathsf T}\bm\alpha-\frac12\bm\alpha^{\mathsf T}YKY\bm\alpha,
 \quad 0\le\alpha_i\le C,\quad \bm y^{\mathsf T}\bm\alpha=0,
 \label{eq:svm-dual}
\end{equation}
其中$\bm\alpha,\bm y\in\R^n$分别收集乘子和标签，$\bm1\in\R^n$是全1向量，$Y=\diag(y_1,\ldots,y_n)$，$K\in\R^{n\times n}$。线性情形$K_{ij}=\bm x_i^{\mathsf T}\bm x_j$。软间隔可以取足够大的正松弛量取得严格可行点，故也满足强对偶。互补条件是
\begin{equation}
 \alpha_i(y_if_i-1+\xi_i)=0,\qquad (C-\alpha_i)\xi_i=0.
\end{equation}
在最优解处，这些条件给出三类系数的含义。$\alpha_i=0$的点满足$y_if_i\ge1$。$0<\alpha_i<C$的点有$\xi_i=0,y_if_i=1$，是自由支持向量。$\alpha_i=C$的点有$y_if_i\le1$，可能在边界上、间隔内或已经错分。

对任一自由支持向量$s$，截距为
\begin{equation}
 b=y_s-\sum_i\alpha_i y_i K_{is}.
 \label{eq:svm-bias}
\end{equation}
实际可对多个自由支持向量取平均。若没有自由支持向量，不能任意选一个$\alpha_s=C$的点套用这个等式。应利用KKT不等式确定$b$的可行区间。令$h_i=\sum_j\alpha_jy_jK_{ji}$。当$\alpha_i=0,y_i=+1$或$\alpha_i=C,y_i=-1$时，有$b\ge y_i-h_i$。另外两种组合给出$b\le y_i-h_i$。对所有下界取最大值、对所有上界取最小值，即得可行区间。数值求解时需为乘子的边界判断设定容差。

\section{特征空间中的非线性分类}
\subsection{特征空间中的样本展开}
若直接在特征空间中求解原始问题，参数$\bm w\in\mathcal H$可能是无限维对象，不能按坐标存储。对偶推导消除了这个参数：由驻点条件，最优方向总在训练特征向量$\phi(\bm x_i)$的张成空间中，因而只需要训练样本之间的内积。把线性情形的内积
\[
 \langle\phi(\bm x_i),\phi(\bm x_j)\rangle_{\mathcal H}
\]
替换成核函数$k(\bm x_i,\bm x_j)$，便得到可计算的特征空间分类器。这个替换必须发生在所有内积位置，包括训练二次项和新样本预测，不能只把最终的线性公式局部替换。

设特征空间中的分类函数为$f(\bm x)=\langle\bm w,\phi(\bm x)\rangle+b$。约束和拉格朗日驻点仍给出
\begin{equation}
 \bm w=\sum_i\alpha_i y_i\phi(\bm x_i),
 \qquad \sum_i\alpha_i y_i=0.
 \label{eq:svm-feature-space-representer}
\end{equation}
代回目标函数后，训练只需要矩阵$K$。因此核方法的非线性来自特征映射，而计算仍在样本对的相似性上进行。

\subsection{核矩阵与预测函数}
若将$\bm x$换成$\phi(\bm x)\in\mathcal H$，上述推导只把内积换成$k(\bm x_i,\bm x_j)$。式\eqref{eq:svm-dual}的结构不变，预测是
\begin{equation}
 f(\bm x)=\sum_{i:\alpha_i>0}\alpha_i y_i k(\bm x_i,\bm x)+b.
\end{equation}
半正定核保证$YKY$半正定，因此对偶是凹最大化问题。预测在特征空间中仍是线性的，在输入空间中可以是曲线。

例如四个点$(\pm1,\pm1)$的标签取$y=x_1x_2$，线性边界不能分开这两类。映射$\phi(\bm x)=(x_1^2,\sqrt2x_1x_2,x_2^2)^{\mathsf T}$使第二个坐标直接区分类别。其内积是$k(\bm x,\bm z)=(\bm x^{\mathsf T}\bm z)^2$。这说明核怎样引入交互项，不说明任意数据在任意核下都可分。相同输入带相反标签时，任何确定性映射都无法分开它们。

\section{求解、验证与适用边界}
对偶有$n$个变量。序列最小优化（Sequential Minimal Optimization，SMO）每次更新少量乘子，同时保持等式约束与箱约束，避免每一步都处理整个二次规划。核缓存可减少重复计算，但大样本下内存仍可能成为限制。多分类通常用一对其余或一对一构造多个二分类器，其分数需要统一决策规则。

训练时先在训练折中拟合缩放统计量，再联合验证$C$和核参数。高斯核长度尺度太小时，不同点的核值接近零，模型容易形成局部响应。尺度太大时，样本相似性接近，模型可能难以区分局部结构。类别不平衡可采用不同的$C_i$或样本权重，但须根据漏报与误报代价验证。

SVM分数不是故障概率。若需要概率输出，可在独立校准数据上拟合分数到概率的映射，再检验校准。最大间隔也不能修复标签泄漏或分布变化。下一章转入无标签数据，研究怎样根据数据的变化寻找低维表示。

\section{阅读依据}
主要依据为Theodoridis第11.10--11.11节（书内页538--552）与MML第12章。郑之杰的核方法文章提供核化公式的对照\cite{zheng2021kernel}，本章补全间隔、对偶与KKT的推导。
